\documentclass[12pt, a4paper, bahasa]{article}

\usepackage[margin=2.5cm]{geometry}
\usepackage{fancyhdr}
\usepackage{hyperref}
\usepackage{graphicx}
\usepackage{float}
\usepackage{titlesec}
\usepackage{wrapfig}
\usepackage{subcaption}
\usepackage{csquotes}
\usepackage{amsmath}
\usepackage{amssymb}
\usepackage{listings}
\usepackage{xcolor}
\usepackage[
    backend=biber,
    style=ieee,
    language=auto,
    autolang=other
]{biblatex}

\addbibresource{refs.bib}

% =========================================================
% HEADER DAN FOOTER
% =========================================================

\pagestyle{fancy}
\setlength{\headheight}{14.5pt}
\fancyfoot[C]{\thepage}
\lhead{\textbf{Nama:} Khalilullah Al Faath}
\rhead{\textbf{NIM:} 25/566643/PPA/07139}
\renewcommand{\headrulewidth}{0.4pt}
\renewcommand{\footrulewidth}{0.4pt}

\renewcommand{\tablename}{Tabel}
\renewcommand{\figurename}{Gambar}

\pagenumbering{arabic}

% =========================================================
% KONFIGURASI LISTING
% =========================================================

\lstset{
    basicstyle=\ttfamily\small,
    frame=single,
    breaklines=true,
    keepspaces=true,
    columns=fullflexible,
    morecomment=[l]{\#},
    commentstyle=\color{red},
    belowskip=0pt,
    showstringspaces=false
}

\begin{document}

\begin{center}
    \large \textbf{Theory of Computation -- Assignment 3} \\[0.2cm]
    \large GOTO Programs \& WHILE Programs \\[0.5cm]
\end{center}

% =========================================================
% REVIEW GOTO PROGRAM
% =========================================================

\section{Review Program GOTO}

Dari slide pada pertemuan ketiga, diketahui bahwa GOTO Program memiliki
ketentuan-ketentuan berikut.
\begin{itemize}
    \item tipe data berupa $\mathbb{N} \in \{0,1,2, ...\}$,
    \item variabel berupa $\{x_0, x_1, x_2, ...\}$,
    \item label berupa $\{L_1, L_2, L_3, ...\}$,
    \item simbol berupa $\{:=, =0, ;, :, +1, -1, \text{IF}, \text{THEN},
    \text{GOTO}, \text{HALT}\}$.
\end{itemize}

Sebuah GOTO program berbentuk daftar baris berlabel
$L_1:S_1;\,L_2:S_2;\,\dots;\,L_n:S_n$, dengan setiap statement $S_i$
berupa salah satu dari lima bentuk berikut (berbeda dengan WHILE/LOOP,
sintaks GOTO \textbf{tidak rekursif}, murni berupa daftar linear).

\begin{enumerate}
    \item $x := x+1$
    \item $x := x-1$
    \item $\texttt{GOTO } L$
    \item $\texttt{IF } x=0 \texttt{ THEN GOTO } L$
    \item $\texttt{HALT}$
\end{enumerate}

\subsection{Semantik}

Instruksi dieksekusi \textbf{berurutan} sesuai urutan label $L_1, L_2,
\dots$, kecuali jika bertemu $\texttt{GOTO } L$ (lompat langsung ke
label $L$) atau $\texttt{IF } x=0 \texttt{ THEN GOTO } L$ (lompat ke
$L$ hanya jika $x=0$, jika tidak lanjut ke baris berikutnya).
Operasi $x:=x+1$/$x:=x-1$ bekerja sama seperti pada WHILE program
(\textit{decrement} tidak menghasilkan bilangan negatif). \texttt{HALT}
menghentikan eksekusi program. Variabel $x_1,\dots,x_k$ diisi input,
variabel lain dimulai dari $0$, dan $x_0$ adalah output.

\subsection{Fungsi GOTO-computable}

Fungsi parsial $f:\mathbb{N}^k \rightharpoonup \mathbb{N}$ disebut
\textbf{GOTO-computable} jika ada GOTO program $P$ berinput $k$
sehingga: jika $f$ terdefinisi pada input $\vec{a}$, maka $P$ berhenti
dengan $x_0 = f(\vec{a})$; jika $f$ tidak terdefinisi pada $\vec{a}$,
maka $P$ tidak pernah berhenti untuk input tersebut.

\subsection{Skema Umum Penerjemahan (dipakai pada Problem 3.2 dan 3.3)}
\label{sec:skema}

Dari pertemuan ketiga, WHILE-computable dan GOTO-computable adalah
himpunan fungsi yang \textbf{identik} (Teorema 3.3 dan Teorema 3.4).
Dua skema penerjemahan berikut dipakai berulang kali pada laporan ini.

\paragraph{Skema GOTO $\to$ WHILE (Teorema 3.3).}
GOTO program $L_1:S_1;\dots;L_n:S_n$ diterjemahkan menjadi satu WHILE
loop besar dengan variabel kontrol $y$ (berperan sebagai \textit{program
counter}):
\begin{lstlisting}
y := 1;
WHILE y != 0 DO
    IF y = 1 THEN P1 END;
    IF y = 2 THEN P2 END;
    ...
    IF y = n THEN Pn END
END
\end{lstlisting}
dengan $P_i$ diturunkan dari $S_i$ menurut aturan:
\begin{itemize}
    \item $S_i = (x:=x\pm1)$ $\Rightarrow$ $P_i = (x:=x\pm1;\ y:=y+1)$,
    \item $S_i = (\texttt{GOTO } L_j)$ $\Rightarrow$ $P_i = (y:=j)$,
    \item $S_i = (\texttt{IF } x=0 \texttt{ THEN GOTO } L_j)$
    $\Rightarrow$ $P_i = (\texttt{IF } x=0 \texttt{ THEN } y:=j
    \texttt{ ELSE } y:=y+1 \texttt{ END})$,
    \item $S_i = \texttt{HALT}$ $\Rightarrow$ $P_i = (y:=0)$.
\end{itemize}
Variabel kontrol $y$ persis meniru nilai label yang sedang aktif
dieksekusi pada GOTO program aslinya; \texttt{WHILE y != 0} berhenti
tepat ketika $y$ diset ke $0$ oleh terjemahan \texttt{HALT}.

\paragraph{Skema WHILE $\to$ GOTO (Teorema 3.4).}
WHILE program $P_1;P_2;\dots;P_n$ diterjemahkan secara rekursif
mengikuti struktur sintaks WHILE, menjadi
$L_1:S_1;\dots;L_m:\texttt{HALT}$, dengan aturan:
\begin{itemize}
    \item $P_i = (x:=x\pm1)$ $\Rightarrow$ langsung jadi satu baris
    $L_j: x:=x\pm1$,
    \item $P_i = (\texttt{WHILE } x\neq0 \texttt{ DO } Q \texttt{ END})$
    $\Rightarrow$ diterjemahkan menjadi
    \begin{lstlisting}
Lj:     IF x = 0 THEN GOTO Ll
        [terjemahan rekursif dari Q]
L(l-1): GOTO Lj
Ll:     [lanjutan setelah loop]
    \end{lstlisting}
\end{itemize}
Setiap $P_i$ bisa menghasilkan lebih dari satu baris GOTO (khususnya
jika $P_i$ berupa WHILE loop, termasuk WHILE loop bersarang), sehingga
label $L_j$ tidak selalu sama dengan indeks $i$.

% =========================================================
% PROBLEM 3.1
% =========================================================

\section{Problem 3.1 -- \textit{infloop}}

Diberikan fungsi parsial
\[
\text{infloop}:\mathbb{N}\rightharpoonup\mathbb{N}, \qquad
\text{untuk semua } x_1\in\mathbb{N},\ \text{infloop}(x_1)\uparrow.
\]
Akan dibuktikan bahwa \textit{infloop} merupakan fungsi GOTO-computable.

\subsection{Ide Penyelesaian}

Karena \textit{infloop} tidak terdefinisi untuk \textbf{semua} input,
berdasarkan definisi GOTO-computable, program yang dikonstruksi harus
\textbf{tidak pernah berhenti}, apapun nilai $x_1$-nya. Cukup dibuat
dua baris yang saling melompat tanpa syarat, tanpa pernah menyentuh
statement \texttt{HALT} dan tanpa pernah melewati baris terakhir
program.

\subsection{Program GOTO}

\begin{lstlisting}
L1: x1 := x1 + 1;
L2: GOTO L1
\end{lstlisting}

Program ini berinput $k=1$ (yaitu $x_1$) dengan $m=1$ (variabel
tertinggi yang dipakai hanya $x_1$).

\subsection{Bukti \textit{Correctness} dan Terminasi}

Akan ditunjukkan bahwa untuk sembarang $x_1\in\mathbb{N}$, program di
atas \textbf{tidak pernah berhenti}, sesuai dengan definisi
GOTO-computable untuk fungsi yang tidak terdefinisi di seluruh domain.

\paragraph{Klaim.} Setelah sejumlah langkah eksekusi berapa pun
(berhingga), label yang sedang aktif selalu berupa $L_1$ atau $L_2$.

\paragraph{Bukti klaim (induksi pada banyaknya langkah).}
\textbf{Basis:} pada langkah ke-0 (sebelum eksekusi dimulai), label
aktif adalah $L_1$, sesuai konvensi bahwa eksekusi GOTO program selalu
dimulai dari $L_1$.
\textbf{Langkah induktif:} andaikan pada suatu langkah label aktif
adalah $L_1$ atau $L_2$.
\begin{itemize}
    \item Jika label aktif $L_1$ (statement $x_1:=x_1+1$, bukan
    lompatan bersyarat maupun \texttt{HALT}), setelah dieksekusi,
    kontrol otomatis lanjut ke baris berikutnya, yaitu $L_2$.
    \item Jika label aktif $L_2$ (statement $\texttt{GOTO } L_1$),
    setelah dieksekusi, kontrol berpindah ke $L_1$.
\end{itemize}
Pada kedua kasus, label aktif berikutnya tetap berupa $L_1$ atau
$L_2$, sehingga klaim terbukti benar untuk semua langkah berikutnya.
$\blacksquare$

\paragraph{Kesimpulan.}
Karena label aktif tidak pernah berupa label lain (khususnya tidak
pernah berupa \texttt{HALT}, sebab tidak ada statement \texttt{HALT}
pada program ini sama sekali), dan program juga tidak pernah melewati
baris terakhir ($L_2$ selalu melompat kembali ke $L_1$, bukan
dilanjutkan ke ``baris sesudah $L_2$'' yang memang tidak ada), maka
eksekusi program \textbf{tidak pernah berhenti}, untuk nilai awal
$x_1$ berapa pun. Hal ini persis sesuai dengan definisi
$\text{infloop}(x_1)\uparrow$ untuk semua $x_1\in\mathbb{N}$. Oleh
karena itu, \textit{infloop} terbukti GOTO-computable. $\blacksquare$

% =========================================================
% PROBLEM 3.2
% =========================================================

\section{Problem 3.2 -- GOTO $\to$ WHILE}

Diberikan fungsi parsial $f:\mathbb{N}^2\rightharpoonup\mathbb{N}$
yang dihitung oleh GOTO program berikut.

\begin{lstlisting}
L1: IF x1 = 0 THEN GOTO L9;
L2: x0 := x0 + 1;
L3: IF x2 = 0 THEN GOTO L9;
L4: x0 := x0 + 1;
L5: x2 := x2 - 1;
L6: x1 := x1 - 1;
L7: GOTO L1;
L8: GOTO L2;
L9: HALT
\end{lstlisting}

Program ini berinput $k=2$ (yaitu $x_1,x_2$), dengan $m=2$ (variabel
tertinggi yang dipakai adalah $x_2$), dan $n=9$ baris. Perhatikan
bahwa baris $L_8$ (\texttt{GOTO L2}) tidak pernah dituju oleh
statement manapun, sehingga merupakan \textit{dead code}; baris ini
tetap harus ikut diterjemahkan agar penerjemahan berlaku umum
(mengikuti skema pada Bagian~\ref{sec:skema}), meskipun tidak
memengaruhi hasil akhir.

\subsection{Ide Penyelesaian}

Sebelum menerjemahkan, akan ditelusuri dahulu apa yang sebenarnya
dihitung oleh program ini, untuk memvalidasi hasil akhir nanti.

Selama $x_1\neq0$ dan $x_2\neq0$ pada awal satu putaran ($L_1$),
program mengeksekusi $L_2$ ($x_0{+}{=}1$), $L_3$ (tidak lompat karena
$x_2\neq0$), $L_4$ ($x_0{+}{=}1$ lagi), $L_5$ ($x_2{-}{=}1$), $L_6$
($x_1{-}{=}1$), lalu kembali ke $L_1$ -- yaitu \textbf{satu putaran
penuh menambah $x_0$ sebanyak $2$, sekaligus mengurangi $x_1$ dan
$x_2$ masing-masing $1$}. Putaran penuh ini berulang sebanyak
$m=\min(x_1,x_2)$ kali, sampai salah satu dari $x_1$ atau $x_2$
mencapai $0$.

Setelah $m$ putaran penuh, ada dua kemungkinan:
\begin{itemize}
    \item Jika $x_1\leq x_2$ (sehingga $x_1$ mencapai $0$ lebih dulu
    atau bersamaan), pada putaran berikutnya $L_1$ langsung mendeteksi
    $x_1=0$ dan lompat ke $L_9$ (\texttt{HALT}) \textbf{tanpa}
    tambahan pada $x_0$. Total: $x_0 = 2m = 2x_1$.
    \item Jika $x_1 > x_2$ (sehingga $x_2$ mencapai $0$ lebih dulu),
    pada putaran berikutnya $L_1$ tidak lompat (karena $x_1\neq0$),
    $L_2$ menambah $x_0$ satu kali lagi, lalu $L_3$ mendeteksi
    $x_2=0$ dan lompat ke $L_9$. Total: $x_0 = 2m+1 = 2x_2+1$.
\end{itemize}

Sehingga secara ringkas,
\[
f(x_1,x_2) =
\begin{cases}
2x_1, & \text{jika } x_1 \leq x_2,\\
2x_2+1, & \text{jika } x_1 > x_2,
\end{cases}
\qquad\text{yaitu}\qquad
f(x_1,x_2) = 2\min(x_1,x_2) +
\begin{cases}
1, & x_1>x_2\\
0, & x_1\leq x_2.
\end{cases}
\]
$f$ adalah fungsi \textbf{total} (program selalu berhenti untuk semua
$x_1,x_2\in\mathbb{N}$, dibuktikan formal pada
Subbagian~\ref{sec:3.2-proof}).

Untuk mengkonstruksi WHILE program yang menghitung $f$, dipakai skema
umum GOTO $\to$ WHILE pada Bagian~\ref{sec:skema}, dengan variabel
kontrol $x_3$ (fresh, belum dipakai di program asal) berperan sebagai
$y$. Karena skema ini terbukti benar secara umum (Teorema 3.3),
WHILE program hasil terjemahan otomatis menghitung $f$ yang sama,
tanpa perlu didesain ulang dari nol.

\subsection{Pemetaan Statement ke $P_i$}

\begin{center}
\begin{tabular}{c|l|l}
$i$ & $S_i$ (GOTO) & $P_i$ (WHILE) \\
\hline
1 & \texttt{IF x1=0 THEN GOTO L9} & \texttt{IF x1=0 THEN x3:=9 ELSE x3:=2 END} \\
2 & \texttt{x0:=x0+1} & \texttt{x0:=x0+1; x3:=3} \\
3 & \texttt{IF x2=0 THEN GOTO L9} & \texttt{IF x2=0 THEN x3:=9 ELSE x3:=4 END} \\
4 & \texttt{x0:=x0+1} & \texttt{x0:=x0+1; x3:=5} \\
5 & \texttt{x2:=x2-1} & \texttt{x2:=x2-1; x3:=6} \\
6 & \texttt{x1:=x1-1} & \texttt{x1:=x1-1; x3:=7} \\
7 & \texttt{GOTO L1} & \texttt{x3:=1} \\
8 & \texttt{GOTO L2} & \texttt{x3:=2} \\
9 & \texttt{HALT} & \texttt{x3:=0} \\
\end{tabular}
\end{center}

\subsection{Program WHILE (hasil terjemahan)}

\begin{lstlisting}
x3 := 1;
WHILE x3 != 0 DO
    IF x3 = 1 THEN
        IF x1 = 0 THEN x3 := 9 ELSE x3 := 2 END
    END;
    IF x3 = 2 THEN
        x0 := x0 + 1;
        x3 := 3
    END;
    IF x3 = 3 THEN
        IF x2 = 0 THEN x3 := 9 ELSE x3 := 4 END
    END;
    IF x3 = 4 THEN
        x0 := x0 + 1;
        x3 := 5
    END;
    IF x3 = 5 THEN
        x2 := x2 - 1;
        x3 := 6
    END;
    IF x3 = 6 THEN
        x1 := x1 - 1;
        x3 := 7
    END;
    IF x3 = 7 THEN
        x3 := 1
    END;
    IF x3 = 8 THEN
        x3 := 2
    END;
    IF x3 = 9 THEN
        x3 := 0
    END
END
\end{lstlisting}

Program ini memakai \textit{syntax sugar} \texttt{IF ... THEN ... END}
dan \texttt{IF ... THEN ... ELSE ... END} yang sudah didefinisikan
pada pertemuan kedua dan ketiga (sesuai catatan soal, \textit{syntax
sugar} diperbolehkan untuk WHILE program).

\subsection{Implementasi Ruby}

\begin{lstlisting}[language=Ruby]
# f.rb
# Implementasi Ruby dari WHILE program hasil terjemahan mekanis
# (skema Teorema 3.3) dari GOTO program f pada Problem 3.2.
# Variabel x3 berperan sebagai variabel kontrol "y" (program counter).

x1 = ARGV[0].to_i
x2 = ARGV[1].to_i

x0 = 0
x3 = 0

# x3 := 1
while x3 != 0
  x3 = x3 - 1
end
x3 = x3 + 1

while x3 != 0

  if x3 == 1
    if x1 == 0
      x3 = 9
    else
      x3 = 2
    end
  end

  if x3 == 2
    x0 = x0 + 1
    x3 = 3
  end

  if x3 == 3
    if x2 == 0
      x3 = 9
    else
      x3 = 4
    end
  end

  if x3 == 4
    x0 = x0 + 1
    x3 = 5
  end

  if x3 == 5
    x2 = x2 - 1
    x3 = 6
  end

  if x3 == 6
    x1 = x1 - 1
    x3 = 7
  end

  if x3 == 7
    x3 = 1
  end

  if x3 == 8   # baris mati (L8 tidak pernah dituju)
    x3 = 2
  end

  if x3 == 9
    x3 = 0
  end

end

puts "f(#{ARGV[0]}, #{ARGV[1]}) = #{x0}"
\end{lstlisting}

\subsection{Hasil Eksekusi}

Program diuji dengan beberapa pasangan input, mencakup kasus
$x_1\leq x_2$, $x_1>x_2$, dan $x_1=x_2$, dibandingkan dengan formula
$f(x_1,x_2)$ hasil analisis semantik di atas.

\begin{lstlisting}
$ ruby source/f.rb 0 0
f(0, 0) = 0
$ ruby source/f.rb 0 5
f(0, 5) = 0
$ ruby source/f.rb 5 0
f(5, 0) = 1
$ ruby source/f.rb 3 0
f(3, 0) = 1
$ ruby source/f.rb 1 3
f(1, 3) = 2
$ ruby source/f.rb 2 1
f(2, 1) = 3
$ ruby source/f.rb 4 4
f(4, 4) = 8
$ ruby source/f.rb 5 2
f(5, 2) = 5
$ ruby source/f.rb 7 3
f(7, 3) = 7
\end{lstlisting}

Seluruh hasil sesuai dengan $f(x_1,x_2)=2\min(x_1,x_2)+[x_1>x_2]$,
misalnya $f(5,2)=2\cdot2+1=5$ dan $f(4,4)=2\cdot4+0=8$.

\subsection{Bukti \textit{Correctness} dan Terminasi}
\label{sec:3.2-proof}

\paragraph{Langkah 1 -- GOTO program asli selalu berhenti.}
Perhatikan bahwa setiap kali kontrol kembali ke $L_1$ dengan
$x_1\neq0$ dan $x_2\neq0$, nilai $\min(x_1,x_2)$ berkurang tepat $1$
pada akhir putaran (karena $x_1$ dan $x_2$ sama-sama dikurangi $1$
pada $L_5$-$L_6$). Karena $\min(x_1,x_2)\in\mathbb{N}$ tidak mungkin
berkurang tanpa batas, setelah paling banyak $\min(x_1,x_2)+1$ putaran
salah satu dari $x_1=0$ (dideteksi di $L_1$) atau $x_2=0$ (dideteksi
di $L_3$) pasti tercapai, sehingga program melompat ke $L_9$ dan
berhenti. Jadi GOTO program berhenti untuk semua $(x_1,x_2)\in
\mathbb{N}^2$, dengan nilai akhir $x_0=f(x_1,x_2)$ seperti diturunkan
pada Ide Penyelesaian.

\paragraph{Langkah 2 -- WHILE hasil terjemahan mensimulasikan GOTO
langkah demi langkah.}
Diklaim: pada setiap saat, nilai variabel kontrol $x_3$ pada WHILE
program sama dengan indeks label yang sedang/akan dieksekusi pada
GOTO program, dan nilai $x_0,x_1,x_2$ pada kedua program selalu sama.

\textit{Basis:} sebelum WHILE loop utama dimulai, $x_3=1$ (dari
$x_3:=1$ di awal), sama seperti GOTO program yang selalu mulai
dieksekusi dari $L_1$; nilai $x_0,x_1,x_2$ sama-sama berupa input
awal.

\textit{Langkah induktif:} andaikan sebelum satu langkah GOTO
$S_i$ dieksekusi (dengan kontrol berada di $L_i$), invarian di atas
berlaku ($x_3=i$, nilai variabel sama). Pada WHILE program, karena
$x_3=i$, blok \texttt{IF x3=i THEN} $P_i$ \texttt{END} adalah
satu-satunya blok pada barisan $P_1,\dots,P_9$ yang kondisinya bisa
bernilai benar terhadap nilai $x_3$ \textbf{sebelum} $P_i$ dieksekusi
(karena kesembilan kondisi $x_3=1,\dots,x_3=9$ saling lepas), sehingga
$P_i$ pasti dieksekusi. Berdasarkan tabel pemetaan, $P_i$ melakukan
persis operasi yang sama pada $x_0,x_1,x_2$ seperti $S_i$ (jika $S_i$
mengubah suatu variabel), dan mengubah $x_3$ menjadi indeks label
berikutnya yang \textbf{persis sama} dengan label berikutnya yang
dituju GOTO program setelah $S_i$ (baik itu $i+1$ untuk statement
biasa, atau $j$ untuk \texttt{GOTO}/\texttt{IF}). Blok-blok
\texttt{IF x3=i'} lain untuk $i' \neq$ nilai baru $x_3$ tidak
mengubah apapun. Sehingga invarian tetap terjaga untuk label
berikutnya. Dengan induksi ini berlaku untuk setiap langkah eksekusi
GOTO program.

\textit{Terminasi bersama:} GOTO program berhenti ketika mengeksekusi
$S_9=\texttt{HALT}$ (Langkah 1), yang oleh invarian bersesuaian dengan
$x_3=9$, sehingga WHILE program mengeksekusi $P_9=(x_3:=0)$, membuat
kondisi \texttt{WHILE x3 != 0} menjadi salah dan loop utama berhenti.
Karena Langkah 1 menjamin GOTO program berhenti untuk semua input,
WHILE program hasil terjemahan juga berhenti untuk semua input, dan
berdasarkan invarian di atas, nilai akhir $x_0$ pada WHILE program
sama dengan nilai akhir $x_0$ pada GOTO program, yaitu $f(x_1,x_2)$.

\paragraph{Kesimpulan.}
WHILE program hasil terjemahan berhenti untuk semua
$(x_1,x_2)\in\mathbb{N}^2$ dan menghasilkan $x_0=f(x_1,x_2)$,
sehingga terbukti WHILE program tersebut menghitung $f$. $\blacksquare$

% =========================================================
% PROBLEM 3.3
% =========================================================

\section{Problem 3.3 -- WHILE $\to$ GOTO}

Diberikan fungsi parsial $g:\mathbb{N}^2\rightharpoonup\mathbb{N}$
yang dihitung oleh WHILE program berikut.

\begin{lstlisting}
WHILE x1 != 0 DO
    WHILE x2 != 0 DO
        x0 := x0 + 1;
        x2 := x2 - 1
    END;
    x1 := x1 - 1
END
\end{lstlisting}

\subsection{Ide Penyelesaian}

Ditelusuri dahulu makna semantik program ini. Selama $x_1\neq0$, pada
setiap putaran loop luar, loop dalam (\texttt{WHILE x2 != 0 DO ...
END}) mengurangi $x_2$ sampai $0$ sambil memindahkan seluruh nilainya
ke $x_0$. Karena $x_2$ tidak pernah diisi ulang di dalam loop luar,
proses ini \textbf{hanya efektif satu kali}, yaitu pada putaran
pertama loop luar (dengan asumsi $x_1>0$): seluruh nilai awal $x_2$
dipindahkan ke $x_0$, dan pada putaran-putaran loop luar berikutnya
loop dalam tidak lagi berbuat apa-apa (karena $x_2$ sudah $0$), loop
luar tersebut hanya mengurangi $x_1$ sampai habis.

Jika $x_1=0$ sejak awal, loop luar tidak pernah dieksekusi sama
sekali, sehingga $x_0$ tetap $0$ (nilai awal), berapapun nilai $x_2$.

Sehingga,
\[
g(x_1,x_2) =
\begin{cases}
x_2, & \text{jika } x_1>0,\\
0, & \text{jika } x_1=0.
\end{cases}
\]
$g$ adalah fungsi \textbf{total} (kedua loop dijamin berhenti, karena
$x_2$ dan $x_1$ hanya berkurang; dibuktikan formal pada
Subbagian~\ref{sec:3.3-proof}).

Untuk mengkonstruksi GOTO program yang menghitung $g$, dipakai skema
umum WHILE $\to$ GOTO pada Bagian~\ref{sec:skema} secara
\textbf{rekursif}: program tersusun dari \texttt{WHILE x1 != 0 DO
$Q$ END}, dengan $Q = (\texttt{WHILE x2 != 0 DO } R \texttt{ END}\,;\,
x_1{:=}x_1{-}1)$, dan $R = (x_0{:=}x_0{+}1\,;\,x_2{:=}x_2{-}1)$.
Loop dalam diterjemahkan lebih dulu, kemudian hasilnya disisipkan ke
dalam terjemahan loop luar.

\subsection{Program GOTO (hasil terjemahan)}

\begin{lstlisting}
L1: IF x1 = 0 THEN GOTO L8;
L2: IF x2 = 0 THEN GOTO L6;
L3: x0 := x0 + 1;
L4: x2 := x2 - 1;
L5: GOTO L2;
L6: x1 := x1 - 1;
L7: GOTO L1;
L8: HALT
\end{lstlisting}

Program ini berinput $k=2$ (yaitu $x_1,x_2$), dengan $m=2$ dan $n=8$
baris. Catatan: sesuai instruksi soal, GOTO program ini ditulis
manual (tidak dibantu Ruby); listing di atas disertakan pula di
laporan digital ini untuk kejelasan, mengikuti versi tulisan tangan.

\paragraph{Korespondensi label.} $L_1$-$L_2$ berperan sebagai
pengujian kondisi loop luar; $L_2$-$L_5$ adalah hasil terjemahan
langsung dari loop dalam \texttt{WHILE x2 != 0 DO $R$ END} (dengan
$L_6$ sebagai label ``setelah loop dalam''); $L_6$ adalah sisa isi
loop luar setelah loop dalam ($x_1:=x_1-1$); $L_7$ melompat kembali
ke pengujian loop luar; $L_8$ adalah label ``setelah loop luar'',
yaitu \texttt{HALT}.

\subsection{Verifikasi dengan Penelusuran}

Sebagai pengecekan tambahan (di luar bukti formal pada
Subbagian~\ref{sec:3.3-proof}), berikut penelusuran untuk
$x_1=3,x_2=5$ (kasus $x_1>0$):

\begin{center}
\begin{tabular}{c|ccc|l}
Label aktif & $x_0$ & $x_1$ & $x_2$ & Keterangan \\
\hline
$L_1$ & 0 & 3 & 5 & $x_1\neq0\to L_2$ \\
$L_2,\dots,L_5$ (5x) & $0\to5$ & 3 & $5\to0$ & loop dalam habiskan $x_2$ \\
$L_2$ & 5 & 3 & 0 & $x_2=0\to L_6$ \\
$L_6$ & 5 & 2 & 0 & $x_1:=x_1-1$ \\
$L_7$ & 5 & 2 & 0 & $\to L_1$ \\
$L_1,L_2,L_6,L_7$ (2x) & 5 & $2\to0$ & 0 & loop dalam tak berbuat apa2 \\
$L_1$ & 5 & 0 & 0 & $x_1=0\to L_8$ \\
$L_8$ & 5 & 0 & 0 & \texttt{HALT} \\
\end{tabular}
\end{center}

Hasil akhir $x_0=5=x_2$ (nilai awal), sesuai $g(3,5)=5$. Untuk
$x_1=0$ (berapapun $x_2$), $L_1$ langsung mendeteksi $x_1=0$ dan
lompat ke $L_8$, sehingga $x_0=0$, sesuai $g(0,x_2)=0$.

\subsection{Bukti \textit{Correctness} dan Terminasi}
\label{sec:3.3-proof}

\paragraph{Langkah 1 -- WHILE program asli selalu berhenti.}
Loop dalam \texttt{WHILE x2 != 0 DO x0:=x0+1; x2:=x2-1 END}
mengurangi $x_2$ tepat $1$ setiap iterasi tanpa pernah menambahnya,
sehingga pasti berhenti setelah tepat nilai $x_2$ saat itu (berhingga)
iterasi. Loop luar \texttt{WHILE x1 != 0 DO ... END} mengurangi $x_1$
tepat $1$ setiap iterasi (pada baris $x_1:=x_1-1$, dieksekusi setelah
loop dalam berhenti), sehingga loop luar juga pasti berhenti, setelah
tepat nilai awal $x_1$ iterasi. Karena kedua loop berhenti (loop dalam
berhenti pada setiap kemunculannya, loop luar berhenti secara
keseluruhan), keseluruhan program berhenti untuk semua
$(x_1,x_2)\in\mathbb{N}^2$, dengan nilai akhir $x_0=g(x_1,x_2)$
seperti diturunkan pada Ide Penyelesaian.

\paragraph{Langkah 2 -- GOTO hasil terjemahan mensimulasikan WHILE
langkah demi langkah.}
Dibuktikan dengan induksi struktural mengikuti struktur sintaks WHILE
program (sesuai skema umum pada Bagian~\ref{sec:skema}):
\begin{itemize}
    \item \textbf{Kasus dasar ($x:=x\pm1$):} diterjemahkan menjadi
    satu baris GOTO yang identik secara semantik, jelas ekivalen.
    \item \textbf{Kasus sequencing ($P_a;P_b$):} jika terjemahan
    $P_a$ dan $P_b$ masing-masing ekivalen (berdasarkan hipotesis
    induksi) dan berhenti dengan kontrol jatuh ke baris pertama
    setelahnya, maka menempatkan terjemahan $P_b$ tepat setelah
    terjemahan $P_a$ menghasilkan program gabungan yang mengeksekusi
    $P_a$ lalu $P_b$ secara berurutan, persis seperti semantik
    $P_a;P_b$. (Diterapkan pada $Q = (\texttt{loop dalam})\,;\,
    (x_1{:=}x_1{-}1)$ untuk menghasilkan $L_2$-$L_6$.)
    \item \textbf{Kasus WHILE loop:} untuk \texttt{WHILE x != 0 DO
    $P$ END} yang diterjemahkan sebagai
    $L_j:\texttt{IF }x{=}0\texttt{ THEN GOTO }L_\ell;\ [\text{terjemahan }
    P];\ L_{\ell-1}:\texttt{GOTO }L_j;\ L_\ell:\dots$ -- berdasarkan
    hipotesis induksi, terjemahan $P$ ekivalen dengan $P$ itu sendiri
    (berhenti dengan efek yang sama, jatuh ke baris pertama
    setelahnya). Maka: jika $x=0$, GOTO langsung lompat ke $L_\ell$
    (loop tidak dieksekusi sama sekali, sama seperti WHILE semantics);
    jika $x\neq0$, GOTO mengeksekusi terjemahan $P$ (satu iterasi),
    lalu $\texttt{GOTO }L_j$ mengulang pengujian $x=0$ kembali --
    persis meniru \texttt{WHILE x != 0 DO P END} yang mengulang
    selama $x\neq0$. Karena bagian ``Langkah 1'' menjamin setiap
    loop pada WHILE program asli berhenti (baik loop dalam maupun
    loop luar), versi GOTO-nya, yang mengeksekusi urutan langkah
    identik, juga pasti mencapai $L_\ell$ (keluar loop) dalam
    berhingga langkah.
\end{itemize}
Diterapkan pada program secara keseluruhan (loop dalam
$\to L_2$-$L_5$, disusul $x_1{:=}x_1{-}1$ di $L_6$, dibungkus loop
luar $\to L_1$-$L_7$, disusul \texttt{HALT} di $L_8$), GOTO program
hasil terjemahan mengeksekusi urutan operasi pada $x_0,x_1,x_2$ yang
identik dengan WHILE program asli pada setiap tahap, dan berhenti
tepat ketika WHILE program asli berhenti.

\paragraph{Kesimpulan.}
GOTO program hasil terjemahan berhenti untuk semua
$(x_1,x_2)\in\mathbb{N}^2$ dan menghasilkan $x_0=g(x_1,x_2)$, sehingga
terbukti GOTO program tersebut menghitung $g$. $\blacksquare$

% =========================================================
% PRANALA KODE
% =========================================================

\section{Pranala Kode}
Hasil \textit{running} implementasi kode dapat diakses di pranala
repositori GitHub berikut: \url{https://github.com/khalilullahalfaath/goto-program-assignment-3-tekkom}.

% =========================================================
% REFERENSI
% =========================================================

% \section{Referensi}
% \printbibliography[heading=none]

\end{document}